% ============================================================
% INTRODUCTION
% ============================================================
\section{Introduction}

Programs that handle medical data sit in an unusual place in the software landscape.
They are usually small and unglamorous---a script that flags abnormal vital signs, a routine that converts a
body weight into a drug dose, a nightly job that exports records to a research registry---yet a mistake in any
of them can reach a patient. Two accidents are often quoted to show how little it takes. The Therac-25 radiation
overdoses of the mid-1980s were traced to defects in control software \cite{leveson1993}, and the Mars Climate
Orbiter was lost in 1999 because one team produced impulse data in pound-force seconds while the consuming
software expected newton seconds \cite{mco1999}. The second failure is, in compiler terms, a plain type error
that the type system of the day simply could not see: both quantities were just \code{double}.

Clinical code has the same weak spot. A dose in milligrams added to a weight in kilograms, a temperature
threshold typed as \code{380} instead of \code{38.0}, a blood-glucose reading that may be missing from the
laboratory feed but is used as if it were always present, or a patient name that reaches a log file without
being anonymised: each of these is syntactically fine and passes through a general-purpose compiler without a
word. The relevant knowledge---units, physiological ranges, privacy rules---lives in the heads of the people who
write the program, not in the language they write it in.

This project builds a small compiler front end that \emph{does} know about that domain, and then asks a
machine-learning question on top of it: once a program has been rejected (or, worse, accepted), can a learned
model tell us what kind of mistake it contains, even when the compiler's own diagnostics are incomplete or
misleading? The generative adversarial network (GAN) enters the picture for a practical reason. Labelled
examples of faulty medical programs are scarce, and they are awkward to share because real clinical scripts tend
to embed identifiers and institutional rules. A conditional GAN that learns the distribution of compiler-derived
feature vectors per error class offers a way to produce additional training records without handing over any
source code. Whether that actually helps is an empirical question, and we report the answer as we measured it,
including the places where it did not help.

\subsection{Motivation}

Three observations motivated the work.

\begin{enumerate}
    \item \textbf{General compilers are blind to domain constraints.} Mainstream languages check that an
    expression is a number, not that it is a mass. Units-of-measure systems exist
    \cite{kennedy1994} but are not part of the languages most clinical scripting is done in.
    \item \textbf{Rule-based diagnostics are precise but narrow.} A hand-written analyser reports exactly what it
    was told to look for. It says nothing about an inverted fever rule such as
    \code{if (temp < 38.0 degC)} followed by an alert called ``fever'', because nothing in the grammar or the
    type system is violated. Cascading faults make things worse: when several defects coexist, the first message
    is not necessarily the root cause.
    \item \textbf{Labelled data is the bottleneck.} Learning-based bug finders depend on large corpora
    \cite{allamanis2018,pradel2018}. In a regulated domain the corpus of real faulty programs is small, and
    synthetic augmentation is attractive if it can be shown to work.
\end{enumerate}

The course context also shaped the design. The assignment asks for a compiler component---here parsing and
semantic analysis---and a machine-learning component---here error classification---applied to a concrete
domain, which for us is healthcare. We therefore treated the compiler as the primary artefact and the learned
model as a second stage that consumes what the compiler produces.

\subsection{Problem Statement}

The problem addressed in this report is:

\begin{quote}
\textit{Given the source text of a medical data-processing program, how can we (i) detect domain-specific
errors---unit, range, privacy and missing-value errors in addition to ordinary lexical, syntactic and type
errors---with a purpose-built compiler front end, and (ii) assign each defective program to its root-cause error
family with a learned classifier that is trained on compiler-derived features, using a generative adversarial
network to compensate for the scarcity of labelled faulty programs?}
\end{quote}

\subsection{Objectives}

\begin{enumerate}
    \item To design \emph{MedScript}, a small imperative language for medical data processing with first-class
    units of measure, protected-health-information (PHI) tagging, optional (possibly missing) values and
    declared physiological ranges.
    \item To implement a complete front end for it: a hand-written lexer, a recursive-descent parser with error
    recovery, and a semantic analyser that performs type, unit, range, taint, missing-value, definite-assignment
    and division-by-zero checks.
    \item To implement the back end far enough to demonstrate the compiler pipeline end to end: three-address
    code, constant folding, constant propagation, common-subexpression elimination, dead-code elimination, and a
    reference virtual machine used to check that optimisation preserves behaviour.
    \item To define a taxonomy of ten classes (valid plus nine error families), build a labelled corpus of valid and fault-injected
    programs, and train classifiers on features extracted from the compiler's tokens, syntax tree and
    diagnostics.
    \item To train conditional GANs on those features and measure honestly whether the synthetic records improve
    classification when real faulty examples are scarce, and whether a GAN trained only on valid programs can
    serve as an anomaly detector.
    \item To measure the compiler itself: scaling of compilation time with program size and the effect of
    optimisation on code size.
\end{enumerate}

\subsection{Contributions}

The main contributions of this report are:

\begin{itemize}
    \item \textbf{MedScript and its compiler.} A domain-specific language and a Python implementation (about
    1{,}900 lines for the compiler proper, about 4{,}900 including corpus tools, learning code, experiments and tests)
    of lexer, parser, semantic analyser, IR generator, optimiser and virtual machine, with a command-line tool that also reports a learned root-cause estimate. A suite of 38 unit tests accompanies it.
    \item \textbf{Semantic analysis that encodes clinical safety rules.} Unit checking through exponent vectors
    over base units, interval reasoning over declared ranges, taint tracking from \code{phi} inputs to output
    sinks, and flow-sensitive refinement of optional values.
    \item \textbf{A fault-injection corpus.} \NProg{} programs in ten classes, produced by 46 mutation operators
    with a check that each injected fault actually manifests, including \DoubleShare\% multi-fault programs among
    the faulty ones.
    \item \textbf{A comparison of rule-based and learned error classification.} The learned classifiers reach
    \RFAcc\% accuracy against \RuleFirstAcc\% for the ``first diagnostic'' rule, and they also catch faults the
    compiler cannot see, which neither rule baseline does.
    \item \textbf{An honest evaluation of GAN augmentation.} Seven training regimes are compared over ten
    random seeds at five levels of data scarcity, together with fidelity measurements of the synthetic records
    and a GAN-based anomaly detector. The GAN turns out not to help, and the report says why.
\end{itemize}

\subsection{Organisation of the Report}

Section~2 reviews related work. Section~3 describes the language, the compiler phases and the machine-learning
components. Section~4 describes the corpus, the experimental protocol and the software environment. Section~5
presents and discusses results. Section~6 concludes and outlines future work. The appendices contain the full
grammar, the diagnostic catalogue, the fault operators and a short guide to running the code.

\newpage
